\part{正文} % 部分标题
% \chapter{章标题}这一章我们介绍这些内容。
\section{章节编排} % 节标题
\subsection{小节标题}这一小节我们介绍这些内容。
\subsubsection{子节标题}这一子节我们介绍这些内容。
\paragraph{段标题}这一段我们介绍这些内容。
\subparagraph{小段标题}这一小段我们介绍这些内容。

此外，在如上命令后添加*，则该部分部分自动编号，且不会编入自动目录中

\section{行文}

\subsection{作者}
\begin{lstlisting}[language={TeX}]
    \usepackage{authblk}
    \title{More than one Author with different Affiliations}
    \author[a]{Author A}
    \author[a]{Author B}
    \author[a]{Author C \thanks{Corresponding author: email@mail.com}}
    \author[b]{Author D}
    \author[b]{Author E}
    \affil[a]{Department of Computer Science, \LaTeX\ University}
    \affil[b]{Department of Mechanical Engineering, \LaTeX\ University}
    % 使用 \thanks 定义通讯作者
    \renewcommand*{\Affilfont}{\small\it} % 修改机构名称的字体与大小
    \renewcommand\Authands{ and } % 去掉 and 前的逗号
    \date{} % 去掉日期
\end{lstlisting}


\subsection{字体}

%字体族设置（罗马字体、无衬线字体、打印机字体）
\textrm{Roman Family} \textsf{Sans Serif Family} \texttt{Typewriter Family}    

{\rmfamily Roman Family} {\sffamily Sans Serif Family} {\ttfamily Typewriter Family}

{\sffamily who you are? you find self on everyone around. take you as the same as others!}

{\ttfamily Are you wiser than others? definitely on. in some days, my it is true. What can you achieve? a luxurious house? a brillilant car? an admirable career? who knows?}

%字体系列设置（粗细、宽度）
\textmd{Medium Series} \textbf{Boldface Series}

{\mdseries Medium Series} {\bfseries Boldface Series}

%字体形状（直立、斜体、伪斜体、小型大写）
\textup{Upright Shape} \textit{Italic Shape} \textsl{Slanted Shape} \textsc{Small Caps Shape}

{\upshape Upright Shape} {\itshape Italic Shape} {\slshape Slanted Shape} {\scshape Small Caps Shape}

%中文字体
{\songti 宋体} \quad {\heiti 黑体} \quad {\fangsong 仿宋} \quad {\kaishu 楷书}

中文字体的 \textbf{粗体} 和 \textit{斜体} 。


%字体大小,根据normalsize的大小确定，normalsize 在文档类的参数决定
{\tiny           Hello}\\
{\scriptsize     Hello}\\
{\footnotesize   Hello}\\
{\small          Hello}\\
{\normalsize      Hello}\\
{\large          Hello}\\
{\Large          Hello}\\
{\LARGE          Hello}\\
{\huge           Hello}\\
{\Huge           Hello}\\

%中文字号设置命令、
\zihao{-0} 你好！
\zihao{5} 你好！

\subsection{目录}
\begin{lstlisting}[language={tex}]
    \phantomsection % 生成新增目录项的链接
    \addcontentsline{toc}{section}{摘要} % 以section的名义向目录中添加条目
\end{lstlisting}

\subsubsection{编号设置}
修改编号的形式
\begin{lstlisting}[language={TeX}]
    \renewcommand\thesection{\arabic{section}} 
    % arabic 阿拉伯数字 
    % roman 小写的罗马数字 
    % Roman 大写的罗马数字 
    % alph 小写字母 
    % Alph 大写字母
    
    同样的，也可以修改子目录的编号方式，而且各目录编号方式可以不同
    \renewcommand\thesection{\Alph{section}} 
    \renewcommand\thesection{\arabic{subsection}}
\end{lstlisting}
其实，Latex内部有23个计数器，17个为序号计数器，6个是控制计数器，选列如下，因此我们可以根据需要改变计数器计数形式。
\begin{table}[h]
    \centering
    \begin{tabular}{|c|c|}
    \hline
    计数器名          & 用途         \\ \hline
    part          & 部序号        \\ \hline
    chapter       & 章序号        \\ \hline
    section       & 节          \\ \hline
    subsection    & 小节         \\ \hline
    subsubsection & 小小节        \\ \hline
    paragraph     & 段          \\ \hline
    subparagraph  & 小段         \\ \hline
    figure        & 插图序号       \\ \hline
    table         & 表格序号       \\ \hline
    equation      & 公式序号       \\ \hline
    page          & 页码计数器      \\ \hline
    footnote      & 脚注序号       \\ \hline
    mpfootnote    & 小页环境中脚注计数器 \\ \hline
    \end{tabular}
    \caption{计数器}
    \label{tab:mytab0}
\end{table}

\subsection{列表}

\paragraph{无序列表}
\begin{itemize}
  \item Javascript
  \item Html
  \item CSS
\end{itemize}

\paragraph{有序列表}
\begin{enumerate}
  \item Javascript
  \item Html
  \item CSS
\end{enumerate}

\paragraph{更加紧凑的无序列表}
\begin{compactitem}
    \item Javascript
    \item Html
    \item CSS
\end{compactitem}

\paragraph{更加紧凑的有序列表}
\begin{compactenum}
    \item Javascript
    \item Html
    \item CSS
\end{compactenum}

\paragraph{描述列表}
\begin{description}
  \item[Javascript] Javascript
  \item[Html] Html
  \item[CSS] CSS
\end{description}


\subsection{换行、分段与分页}
\begin{itemize}
    \item $\backslash \backslash$ 表示强制换行 与 $\backslash newline$ 类似但有不同
    \item $\backslash \backslash[offset]$ 表示强制换行并且与下一行的行间距为原来行间距+offset
    \item $\backslash linebreak$ 强制换行，与$\backslash newline$的区别为其当前行分散对齐。
    \item $\backslash par$ 表示强制分段
    \item $\backslash newpage$ 分页命令。
    \item $\backslash clearpage$ 分页命令
\end{itemize}

\begin{lstlisting}
    \newpage：  The \newpage command ends the current page.
    \clearpage：The \clearpage command ends the current page and causes all figures and tables that have so far appeared in the input to be printed.
\end{lstlisting}

通俗点讲就是当你新加的一页内容较多时，两者基本一样，当新加的一页内容较少时，$\backslash newpage$就无法实现你想要的效果，但$\backslash clearpage$可以。


\subsection{空格}
$\backslash quad, 1em, em, m$ 代表当前字体下接近字符‘M’的宽度。

\begin{table}[h]
    \centering
    \begin{tabular}{|c|c|}
        \hline
        两个quad空格 & a \qquad b \\
        \hline
        一个quad空格 & a \quad b \\
        \hline
        大空格	 & a\ b \\
        \hline
        中等空格 &	a\;b \\
        \hline
        小空格 & a\,b \\
        \hline
        没有空格 & ab \\
        \hline
        紧贴    &  a\!b  \\
        \hline
    \end{tabular}
    \caption{空格效果}
    \label{tab:mytab1}
\end{table}

\subsection{缩进}
$\backslash noindent$ 可以强制首行不缩进

$\backslash setlength \{\backslash parindent\}\{1em\}$ 自行设置缩进量

\subsection{引用}
在bib文件中写入数据，并以$\backslash cite{}$引用，在文档最后呈现所有参考文献。

引用共三种 $cite,url/href,ref$ 分别是
\begin{itemize}
    \item 对bib中的引用 \cite{W:liam} \cite{W:jingwhale}
    \item 对网址的直接引用 \url{www.baidu.com} / \href{www.baidu.com}{百度}
    \item 对label的引用 如图\ref{fig:myfig1}
\end{itemize}

此外，可以使用$\backslash footnote \{脚注内容\}$添加脚注(标题中添加脚注时需要在在前面多加$\backslash protect$ )，脚注标号设置如下
\begin{lstlisting}[language={TeX}]
    % 符号/字母 脚注标号设置
    \renewcommand{\thefootnote}{\fnsymbol{footnote}} % 重置\footnote{}命令为符号脚注
    \renewcommand{\thefootnote}{\alph{footnote}} % 重置\footnote{}命令为符号脚注
    
    % arabic 阿拉伯数字 
    % roman 小写的罗马数字 
    % Roman 大写的罗马数字 
    % alph 小写字母 
    % Alph 大写字母
    
    % 数字脚注标号设置，且标号可以略去
    \usepackage{lipsum}
    \newcommand\blfootnote[1]{% 
    \begingroup 
    \renewcommand\thefootnote{}\footnote{#1}% 
    \addtocounter{footnote}{-1}% 
    \endgroup 
    }
    
    \blfootnote{*Corresponding Author} % 无标号
    \footnote{Corresponding Author} % 正常显示标号，从1开始
\end{lstlisting}


\subsection{代码}
使用listings与xcolor实现代码段的插入，以下是配置示例
\begin{lstlisting}[language={TeX}]
\usepackage{listings} % 代码段
\usepackage{xcolor} % 颜色

\lstset{ %
  backgroundcolor=\color{white},   % choose the background color; you must add \usepackage{color} or \usepackage{xcolor}
  basicstyle=\ttfamily,            % the size of the fonts that are used for the code
  breakatwhitespace=false,         % sets if automatic breaks should only happen at whitespace
  breaklines=true,                 % sets automatic line breaking
  captionpos=b,                    % sets the caption-position to bottom
  commentstyle=\ttfamily\color{mygreen},    
                                   % comment style
  deletekeywords={},               % if you want to delete keywords from the given language
  escapeinside={},                 % if you want to add LaTeX within your code
  extendedchars=true,              % lets you use non-ASCII characters; for 8-bits encodings only, does not work with UTF-8
  frame=single,                    % adds a frame around the code
  keepspaces=true,                 % keeps spaces in text, useful for keeping indentation of code (possibly needs columns=flexible)
  keywordstyle=\color{blue},       % keyword style
  language=C++,                    % the language of the code
  morekeywords={},                 % if you want to add more keywords to the set
  numbers=left,                    % where to put the line-numbers; possible values are (none, left, right)
  numbersep=5pt,                   % how far the line-numbers are from the code
  numberstyle=\tiny\color{mygray}, % the style that is used for the line-numbers
  rulecolor=\color{black},         % if not set, the frame-color may be changed on line-breaks within not-black text (e.g. comments (green here))
  showspaces=false,                % show spaces everywhere adding particular underscores; it overrides 'showstringspaces'
  showstringspaces=false,          % underline spaces within strings only
  showtabs=false,                  % show tabs within strings adding particular underscores
  stepnumber=1,                    % the step between two line-numbers. If it's 1, each line will be numbered
  stringstyle=\color{mymauve},     % string literal style
  tabsize=2,                       % sets default tabsize to 2 spaces
  title=\lstname                   % show the filename of files included with \lstinputlisting; also try caption instead of title
}
\end{lstlisting}

\subsection{算法}

\section{数学公式}

LaTeX 的数学模式有两种：行内模式 (inline) 和行间模式 (display)。前者在正文的行文中，插入数学公式；后者独立排列单独成行，并自动居中。

在行文中，使用 $ X^2 $ 可以插入行内公式，使用 \[ X^2 \] 可以插入行间公式，如果需要对行间公式进行编号，则可以使用 equation 环境（equation表示不需要编号）：

\begin{equation}
\sqrt{x}, \frac{1}{2}
\end{equation}

\subsection{运算符} 

一些小的运算符，可以在数学模式下直接输入；另一些需要用控制序列生成，如
\[ \pm\; \times \; \div\; \cdot\; \cap\; \cup\;
\geq\; \leq\; \neq\; \approx \; \equiv \]

连加、连乘、极限、积分等大型运算符 % \sum, \prod, \lim, \int
他们的上下标在行内公式中被压缩，以适应行高。我们可以强制显式地指定是否压缩这些上下标。例如：

$ \sum_{i=1}^n i\quad \prod_{i=1}^n $ 

$ \sum\limits _{i=1}^n i\quad \prod\limits _{i=1}^n $

\[ \lim_{x\to0}x^2 \quad \int_a^b x^2 dx \]
\[ \lim\nolimits _{x\to0}x^2\quad \int\nolimits_a^b x^2 dx \]

多重积分
\[ \iint\quad \iiint\quad \iiiint\quad \idotsint \]

\subsection{定界符}

\[ \Biggl(\biggl(\Bigl(\bigl((x)\bigr)\Bigr)\biggr)\Biggr) \]
\[ \Biggl[\biggl[\Bigl[\bigl[[x]\bigr]\Bigr]\biggr]\Biggr] \]
\[ \Biggl \{\biggl \{\Bigl \{\bigl \{\{x\}\bigr \}\Bigr \}\biggr \}\Biggr\} \]
\[ \Biggl\langle\biggl\langle\Bigl\langle\bigl\langle\langle x
\rangle\bigr\rangle\Bigr\rangle\biggr\rangle\Biggr\rangle \]
\[ \Biggl\lvert\biggl\lvert\Bigl\lvert\bigl\lvert\lvert x
\rvert\bigr\rvert\Bigr\rvert\biggr\rvert\Biggr\rvert \]
\[ \Biggl\lVert\biggl\lVert\Bigl\lVert\bigl\lVert\lVert x
\rVert\bigr\rVert\Bigr\rVert\biggr\rVert\Biggr\rVert \]

\subsection{省略号}
\[ x_1,x_2,\dots ,x_n\quad 1,2,\cdots ,n\quad
\vdots\quad \ddots \]

\subsection{矩阵}
\[ \begin{pmatrix} a&b\\c&d \end{pmatrix} \quad
\begin{bmatrix} a&b\\c&d \end{bmatrix} \quad
\begin{Bmatrix} a&b\\c&d \end{Bmatrix} \quad
\begin{vmatrix} a&b\\c&d \end{vmatrix} \quad
\begin{Vmatrix} a&b\\c&d \end{Vmatrix} \]

行内的小矩阵：Marry has a little matrix $ ( \begin{smallmatrix} a&b\\c&d \end{smallmatrix} ) $.

\subsection{多行公式}

\subsubsection{长公式}
不对齐
\begin{multline}
x = a+b+c+{} \\
d+e+f+g
\end{multline}
对齐
\[\begin{aligned}
x ={}& a+b+c+{} \\
&d+e+f+g
\end{aligned}\]

\subsubsection{公式组}
\begin{gather}
a = b+c+d \\
x = y+z
\end{gather}
\begin{align}
a &= b+c+d \\
x &= y+z
\end{align}

\subsubsection{分段公式}
\begin{equation}
    y= \begin{cases}
    -x,\quad x\leq 0 \\
    x,\quad x>0
    \end{cases} 
\end{equation}

\subsubsection{复杂公式}
\begin{equation}
\left.
\begin{aligned}
x+y &> 5 \\
y-y &> 11
\end{aligned}
\ \right\}\Rightarrow x^2 - y^2 > 55
\end{equation}


\section{图表}
$\backslash centerline{}$ 以及 $\backslash centering$可以居中图表

\subsection{插入图片}
如图\ref{fig:myfig1}
\begin{figure}[htbp]
\centering
\includegraphics[width = .8\textwidth]{figure/Latex.jpg}
\caption{图片标题}
\label{fig:myfig1}
\end{figure}

\subsection{插入表格}
如表\ref{tab:mytab2}，tabular环境提供了最简单的表格功能。它用 $\backslash hline$ 命令表示横线，在列格式中用 | 表示竖线；用 $\&$ 来分列，用 $\backslash\backslash$换行 

Latex的表格编写实在是太繁琐了，推荐使用 \textbf{\href{http://www.tablesgenerator.com/}{Tables Generator}} 轻松愉悦的完成表格 :>

\begin{table}[htbp]
    \centering
    \begin{tabular}{|l|c|r|}
         \hline
        操作系统& 发行版& 编辑器\\
         \hline
        Windows & MikTeX & TexMakerX \\
         \hline
        Unix/Linux & teTeX & Kile \\
         \hline
        Mac OS & MacTeX & TeXShop \\
         \hline
        通用& TeX Live & TeXworks \\
         \hline
    \end{tabular}
    \caption{表格标题}
    \label{tab:mytab2}
\end{table}

\subsubsection{跨列表格}
如表\ref{tab:mytab3}
\begin{table}[htbp]
    \centering
    \begin{tabular}{|l|c|r|}
      \hline
      % after \\: \hline or \cline{col1-col2} \cline{col3-col4} ...
      左列 & 中列 & 右列 \\
      \hline
      2行1列 & 2行2列 & 2行3列 \\
      \hline
      \multicolumn{2}{|c|}{跨越2015} & 3行3列 \\
      \hline
      4行1列 & 4行2列 & 4行3列 \\
      \hline
    \end{tabular}
    \caption{跨列表格}
    \label{tab:mytab3}
\end{table}

\subsection{浮动效果}
 htbp 选项用来指定插图的理想位置，这几个字母分别代表 here, top, bottom, float page，也就是就这里、页顶、页尾、浮动页（专门放浮动体的单独页面或分栏）。$\backslash centering$ 用来使插图居中；$\backslash caption$ 命令设置插图标题，LaTeX 会自动给浮动体的标题加上编号。注意 $\backslash label$ 应该放在标题命令之后。

\subsection{图表混合}
\subsubsection{一行多图}
在导言区引入以下包，使用subfigure即可实现一行多图等效果，并且图片的标题与label都可以相应设置。子图的标题中的字母可以通过设置修改
\begin{lstlisting}[language={TeX}]
    \usepackage{graphicx}
    \usepackage{caption}
    \usepackage{subcaption}
\end{lstlisting}
$\backslash quad$分开两张图片，$\backslash quad$之后再加上回车，则分行

\begin{figure}[htbp]
    \centering
    \begin{subfigure}[t]{0.2\linewidth}
        \includegraphics[width=\linewidth]{figure/Latex.jpg}
        \caption{一行多图}
        \label{fig:myfig2}
    \end{subfigure}
    \quad
    \begin{subfigure}[t]{0.2\linewidth}
        \includegraphics[width=\linewidth]{figure/Latex.jpg}
        \caption{一行多图}
        \label{fig:myfig3}
    \end{subfigure}
    \quad
    \begin{subfigure}[t]{0.2\linewidth}
        \includegraphics[width=\linewidth]{figure/Latex.jpg}
        \caption{一行多图}
        \label{fig:myfig4}
    \end{subfigure}
    \quad
    \begin{subfigure}[t]{0.2\linewidth}
        \includegraphics[width=\linewidth]{figure/Latex.jpg}
        \caption{一行多图}
        \label{fig:myfig5}
    \end{subfigure}
    \caption{一行多图}
    \label{fig:myfig6}
\end{figure}

\begin{figure}[htbp]
    \centering
    \begin{subfigure}[t]{0.4\linewidth}
        \includegraphics[width=\linewidth]{figure/Latex.jpg}
        \caption{多行多图}
        \label{fig:myfig7}
    \end{subfigure}
    \quad
    \begin{subfigure}[t]{0.4\linewidth}
        \includegraphics[width=\linewidth]{figure/Latex.jpg}
        \caption{多行多图}
        \label{fig:myfig8}
    \end{subfigure}
    \quad
    
    \begin{subfigure}[t]{0.4\linewidth}
        \includegraphics[width=\linewidth]{figure/Latex.jpg}
        \caption{多行多图}
        \label{fig:myfig9}
    \end{subfigure}
    \quad
    \begin{subfigure}[t]{0.4\linewidth}
        \includegraphics[width=\linewidth]{figure/Latex.jpg}
        \caption{多行多图}
        \label{fig:myfig10}
    \end{subfigure}
    \caption{多行多图}
    \label{fig:myfig11}
\end{figure}

\subsubsection{一行多表}
一行多表以及图表混合均使用minipage的方式实现，标题与label均独立，可以在\href{http://www.sascha-frank.com/latex-minipage.html}{这里}查看其文档

\begin{lstlisting}[language={TeX}]
    \begin{minipage}[adjusting]{width of the minipage}
     Text ... \ \
     Images ... \ \
     Tables ... \ \
    \end{minipage} 
\end{lstlisting}

\begin{table}[htbp]
  \begin{minipage}{0.45\linewidth}
    \centering
    \captionof{table}{一行多表}
    \begin{tabular}{|l|c|r|}
         \hline
        操作系统& 发行版& 编辑器\\
         \hline
        Windows & MikTeX & TexMakerX \\
         \hline
        Unix/Linux & teTeX & Kile \\
         \hline
        Mac OS & MacTeX & TeXShop \\
         \hline
        通用& TeX Live & TeXworks \\
         \hline
    \end{tabular}
    \label{tab:mytab4}
  \end{minipage}
  \quad
  \begin{minipage}{0.45\linewidth}
    \centering
    \captionof{table}{一行多表}
        \begin{tabular}{|l|c|r|}
             \hline
            操作系统& 发行版& 编辑器\\
             \hline
            Windows & MikTeX & TexMakerX \\
             \hline
            Unix/Linux & teTeX & Kile \\
             \hline
            Mac OS & MacTeX & TeXShop \\
             \hline
            通用& TeX Live & TeXworks \\
             \hline
        \end{tabular}
    \label{tab:mytab5}
  \end{minipage}
\end{table}

\subsubsection{图表混合}

\begin{figure}[htb]
    \centering
    \begin{minipage}[c]{0.4\linewidth}
        \centering
        \includegraphics[width=\linewidth]{figure/Latex.jpg}
        \caption{图表混合}
        \label{fig:myfig12}
    \end{minipage}
    \quad
    \begin{minipage}[c]{0.4\linewidth}
        \centering
        \begin{tabular}{|l|c|r|}
             \hline
            操作系统& 发行版& 编辑器\\
             \hline
            Windows & MikTeX & TexMakerX \\
             \hline
            Unix/Linux & teTeX & Kile \\
             \hline
            Mac OS & MacTeX & TeXShop \\
             \hline
            通用& TeX Live & TeXworks \\
             \hline
        \end{tabular}
        \captionof{table}{图表混合}
        \label{tab:mytab6}
    \end{minipage}
\end{figure}

\section{版面设置}

\subsection{段落对齐}

\begin{center}
对于生成复杂表格和数学公式，这一点表现得尤为突出。
\end{center}
\begin{flushleft}
LATEX（英语发音：/ˈleɪtɛk/ ），
\end{flushleft}
\begin{flushright}
是一种基于TEX的排版系统，
\end{flushright}

\subsection{页面边距}

设置页边距，推荐使用 geometry 宏包。可以在\textbf{\href{http://texdoc.net/texmf-dist/doc/latex/geometry/geometry.pdf}{这里}}查看它的说明文档。

例如将纸张的长度设置为 20cm、宽度设置为 15cm、左边距 1cm、右边距 2cm、上边距 3cm、下边距 4cm，后两行可以定制特定页的边距：
\begin{lstlisting}[language={TeX}]
    \usepackage{geometry}
    \geometry{papersize={20cm,15cm}}
    \geometry{left=1cm,right=2cm,top=3cm,bottom=4cm}
    \newgeometry{left = 0.8 cm, right = 0.8 cm, bottom = 0.8cm} % 单独设置从此行向下的页边距
    \restoregeometry % 还原页边距
\end{lstlisting}

\subsection{页眉页脚}
设置页眉页脚，推荐使用 fancyhdr 宏包。可以在\textbf{\href{http://texdoc.net/texmf-dist/doc/latex/fancyhdr/fancyhdr.pdf}{这里}}查看它的说明文档。

自定义页眉和页脚。
为页眉和页脚加上装饰性的横线。
页眉和页脚的宽度可以超过正文文本的宽度。
多行的页眉和页脚。
奇偶页使用不同格式的页眉和页脚。
每章的首页使用不同格式的页眉和页脚。
浮动对象页使用不同格式的页眉和页脚。
控制页眉和页脚的字体，包括字形，字族，大小写等。

比如我希望，在页眉左边写上我的名字，中间写上今天的日期，右边写上我的电话；页脚的正中写上页码；页眉和正文之间有一道宽为 0.4pt 的横线分割，可以在导言区加上如下几行：

\begin{lstlisting}[language={TeX}]
    \usepackage{fancyhdr}
    \pagestyle{fancy}
    \lhead{\author}
    \chead{\date}
    \rhead{}
    \lfoot{}
    \cfoot{\thepage}
    \rfoot{}
    \renewcommand{\headrulewidth}{0.4pt}
    \renewcommand{\headwidth}{\textwidth}
    \renewcommand{\footrulewidth}{0pt}
    
    \pagenumbering{数字形式} % arabic, 阿拉伯数字 roman, 小写罗马数字 Roman,大写罗马数字 alpha, 小写拉丁字母 Alpha, 大写拉丁字母
\end{lstlisting}

\subsection{行间距}
我们可以通过 setspace 宏包提供的命令来调整行间距。比如在导言区添加如下内容，可以将行距设置为字号的 1.5 倍：
\begin{lstlisting}[language={TeX}]
    \usepackage{setspace}
    \onehalfspacing
\end{lstlisting}

\subsection{段间距}
我们可以通过修改长度 $\backslash parskip$ 的值来调整段间距。例如在导言区添加以下内容

\begin{lstlisting}[language={TeX}]
    \addtolength{\parskip}{.4em}
\end{lstlisting}

则可以在原有的基础上，增加段间距0.4em。如果需要减小段间距，只需将该数值改为负值即可。

\bibliographystyle{plain} % 指定参考文献的呈现方式
\clearpage
\phantomsection
\addcontentsline{toc}{section}{参考文献} %向目录中添加条目，以章的名义
\bibliography{ref} % 命令用于指定之前生成的.bib库。